\chapter{मूलभूत अंतःक्रियाएँ}\label{ch:g11:fundamental-interactions}

गुब्बारे को बालों पर रगड़िए और वह काग़ज़ के टुकड़े उठा लेगा --- रबर के चंद
वर्ग सेंटीमीटर से पूरे ग्रह के \omterm{def:g11:fundamental-interactions:four}{गुरुत्वाकर्षण} को हराकर। इसी आसान जीत के कारण
\omterm{def:g11:fundamental-interactions:ladder}{परमाणु} जुड़े रहते हैं और भारी नाभिक अंततः टूट जाते हैं (\cref{ch:g11:nucleus-radioactivity})।
ब्रह्मांड ठीक चार अंतःक्रियाओं पर चलता है; यह अध्याय उनसे मिलता है, उन्हें
आपस में तौलता है और हर एक को उसकी मंज़िल सौंप देता है।

\section{क्वार्क से आकाशगंगा तक का पदार्थ}

\begin{definition}[संरचना की सीढ़ी]\label{def:g11:fundamental-interactions:ladder}
पदार्थ मंज़िलों में बना है, और हर मंज़िल नीचे वाली से जुड़कर बनी है: तीन
\emph{क्वार्क}\index{क्वार्क} (बिंदुवत्, \qty{e-18}{m} से भी छोटे) मिलकर एक
\emph{न्यूक्लियॉन}\index{न्यूक्लियॉन} बनाते हैं --- प्रोटॉन या न्यूट्रॉन,
लगभग \qty{e-15}{m} चौड़ा; न्यूक्लियॉन ठुँसकर \emph{नाभिक}\index{नाभिक} बनाते हैं
(\qty{e-15}{m} से \qty{e-14}{m} तक); नाभिक और उसका इलेक्ट्रॉन-बादल मिलकर
\emph{परमाणु}\index{परमाणु} हैं (\qty{e-10}{m}); परमाणु जुड़कर
\emph{अणु}\index{अणु} और क्रिस्टल बनाते हैं (\qty{e-9}{m} से ऊपर), जो रोज़मर्रा की
वस्तुओं (\qty{1}{m}), ग्रहों (\qty{e7}{m}), ग्रह-मंडलों (\qty{e11}{m}), आकाशगंगाओं
(\qty{e21}{m}) और प्रेक्षणीय ब्रह्मांड (\qty{e26}{m}) में जुड़ते जाते हैं।
\end{definition}

\begin{definition}[चार अंतःक्रियाएँ]\label{def:g11:fundamental-interactions:four}
\emph{मूलभूत अंतःक्रिया}\index{मूलभूत अंतःक्रिया} वह बल है जिसे
किन्हीं और बलों में नहीं तोड़ा जा सकता; ऐसी ठीक चार हैं:
\emph{गुरुत्वाकर्षण}\index{गुरुत्वाकर्षण}, \emph{विद्युत्-चुंबकीय
अंतःक्रिया}\index{विद्युत्-चुंबकीय अंतःक्रिया}, \omterm{def:g11:fundamental-interactions:strong}{प्रबल अंतःक्रिया} और
\omterm{def:g11:fundamental-interactions:weak}{दुर्बल अंतःक्रिया}।
\end{definition}

\begin{remark}\label{rem:g11:fundamental-interactions:empty}
मंज़िलों के बीच सूनापन है: \omterm{def:g11:fundamental-interactions:ladder}{परमाणु} अपने नाभिक से $10^{5}$ गुना
चौड़ा है --- नाभिक को बढ़ाकर \qty{1}{cm} की कंचे जितना कर दीजिए तो
इलेक्ट्रॉन-बादल एक किलोमीटर फैल जाएगा। इन बिखरी मंज़िलों को कुछ तो थामे रखता
होगा, और वह गुरुत्व लगभग कभी नहीं होता।
\end{remark}

\begin{omfigure}
\begin{tikzpicture}[font=\small, x=0.24cm]
  \draw[-Stealth] (-1,0) -- (46.5,0);
  \foreach \x/\l/\d in {0/{10^{-18}}/2pt, 3/{10^{-15}}/13pt,
                        8/{10^{-10}}/2pt, 18/{10^{0}}/2pt,
                        25/{10^{7}}/2pt, 29/{10^{11}}/13pt,
                        39/{10^{21}}/2pt, 44/{10^{26}}/13pt}
    \draw (\x,0.1) -- (\x,-0.1) node[below=\d] {$\l$};
  \foreach \x/\n in {0/quark, 3/nucleus, 8/atom, 18/human, 25/Earth,
                     29/{Sun--Earth}, 39/galaxy, 44/universe}{
    \fill[omDef] (\x,0) circle (1.6pt);
    \node[rotate=45, anchor=west, inner sep=1.5pt] at (\x,0.24) {\n};}
\end{tikzpicture}

{\small दस की घातों वाले अक्ष पर संरचना की सीढ़ी: \omterm{def:g11:fundamental-interactions:ladder}{क्वार्क} और
प्रेक्षणीय ब्रह्मांड के बीच चौवालीस दहाइयाँ हैं। आकार \omterm{def:g10:orders-of-magnitude:unit}{मीटर} में।}
\end{omfigure}

\section{विद्युत आवेश}

\begin{definition}[विद्युत आवेश]\label{def:g11:fundamental-interactions:charge}
\emph{विद्युत आवेश}\index{विद्युत आवेश} पदार्थ का वह गुण है जिस पर
\omterm{def:g11:fundamental-interactions:four}{विद्युत्-चुंबकीय अंतःक्रिया} लगती है, ठीक जैसे द्रव्यमान वह गुण है
जिस पर \omterm{def:g11:fundamental-interactions:four}{गुरुत्वाकर्षण} लगता है। इसे \emph{कूलॉम}\index{कूलॉम} (\unit{C})
में मापते हैं और यह \emph{दो चिह्नों} में आता है, धनात्मक और ऋणात्मक, जो
जुड़ने पर एक-दूसरे को काट देते हैं; समान आवेश एक-दूसरे को प्रतिकर्षित करते
हैं, विपरीत आवेश आकर्षित।
\end{definition}

\begin{definition}[मूल आवेश]\label{def:g11:fundamental-interactions:elementary}
आवेश दानेदार है: हर प्रेक्षित आवेश \emph{मूल आवेश}\index{मूल आवेश} $e = \qty{1.60e-19}{C}$ का
पूरा गुणज होता है। प्रोटॉन $+e$ रखता है, इलेक्ट्रॉन $-e$, न्यूट्रॉन
$0$; और आवेश $q$ में $N = \abs{q}/e$ मूल आवेश होते हैं।
\end{definition}

\begin{proposition}[आवेश का संरक्षण]\label{prop:g11:fundamental-interactions:conservation}
किसी विलगित निकाय का कुल आवेश कभी नहीं बदलता: आवेश न बनता है न नष्ट होता है,
केवल हस्तांतरित होता है --- व्यवहार में इलेक्ट्रॉनों के ज़रिए, जो सबसे हल्के
वाहक हैं।
\end{proposition}
\admitted

\begin{remark}\label{rem:g11:fundamental-interactions:conservation}
इसका कोई उल्लंघन कभी नहीं देखा गया; इसके पीछे का गहरा कारण, विद्युत्-चुंबकत्व
की एक सममिति, विश्वविद्यालय के खंडों में दिया गया है।
\end{remark}

\begin{definition}[घर्षण और संपर्क से आवेशन]\label{def:g11:fundamental-interactions:charging}
दो विद्युत्रोधियों को रगड़ने पर इलेक्ट्रॉन एक से उखड़कर दूसरे पर चले जाते हैं
और दोनों पर बराबर परिमाण के विपरीत आवेश छूट जाते हैं (\emph{घर्षण द्वारा
आवेशन}\index{घर्षण द्वारा आवेशन}); किसी आवेशित वस्तु को उदासीन वस्तु से छुआने
पर आवेश बँट जाता है (\emph{संपर्क द्वारा आवेशन}\index{संपर्क द्वारा आवेशन})।
हिलते केवल इलेक्ट्रॉन हैं; कुल आवेश संरक्षित रहता है।
\end{definition}

\begin{example}[इलेक्ट्रॉन गिनना]\label{ex:g11:fundamental-interactions:balloon}
बालों पर रगड़े गुब्बारे को $q = \qty{-1.6e-8}{C}$ मिलता है: उसने
$N = \num{1.6e-8} / \num{1.6e-19} = \num{1.0e11}$ इलेक्ट्रॉन \emph{पाए} हैं, और बालों पर ठीक $+\qty{1.6e-8}{C}$ बच रहता है।
\end{example}

\section{कूलॉम का नियम}

\begin{theorem}[कूलॉम का नियम]\label{thm:g11:fundamental-interactions:coulomb}
$d$ दूरी पर रखे दो बिंदु-आवेश $q_1$ और $q_2$ एक-दूसरे पर उन्हें जोड़ती
रेखा के अनुदिश बल लगाते हैं, जिनके परिमाण बराबर और दिशाएँ विपरीत
होती हैं --- समान चिह्नों पर प्रतिकर्षी, विपरीत चिह्नों पर आकर्षी --- और
\[
F = k\,\frac{\abs{q_1 q_2}}{d^{2}},
\qquad k = \qty{8.99e9}{N.m^2/C^2}.
\]
\end{theorem}
\admitted

\begin{remark}\label{rem:g11:fundamental-interactions:torsion}
\omterm{def:g11:fundamental-interactions:charge}{कूलॉम} ने यह नियम 1785 में एक \omterm{def:g11:lenses-and-eye:lens}{पतले} ऐंठन-तंतु का मरोड़ नापकर स्थापित
किया था; \omterm{def:g10:orders-of-magnitude:scinot}{घातांक} ठीक $2$ ही क्यों है, इसका उत्तर विश्वविद्यालय के
खंडों में है।
\end{remark}

\begin{omfigure}
\begin{tikzpicture}[font=\small]
  \draw[omDef, thick, fill=omDef!20] (0,0) circle (0.3) node[black] {$+$};
  \draw[omDef, thick, fill=omDef!20] (2.6,0) circle (0.3) node[black] {$+$};
  \draw[-Stealth, omThm, very thick] (-0.3,0) -- (-1.5,0)
    node[midway, above] {$\vect F_{2/1}$};
  \draw[-Stealth, omThm, very thick] (2.9,0) -- (4.1,0)
    node[midway, above] {$\vect F_{1/2}$};
\end{tikzpicture}\qquad
\begin{tikzpicture}[font=\small]
  \draw[omDef, thick, fill=omDef!20] (0,0) circle (0.3) node[black] {$+$};
  \draw[omProp, thick, fill=omProp!25] (2.6,0) circle (0.3) node[black] {$-$};
  \draw[-Stealth, omThm, very thick] (0.3,0) -- (1.1,0)
    node[midway, above] {$\vect F_{2/1}$};
  \draw[-Stealth, omThm, very thick] (2.3,0) -- (1.5,0)
    node[midway, above] {$\vect F_{1/2}$};
\end{tikzpicture}

{\small समान आवेश प्रतिकर्षित करते हैं (बाएँ), विपरीत आवेश आकर्षित (दाएँ);
दोनों बलों के परिमाण हमेशा बराबर और दिशाएँ विपरीत रहती हैं।}
\end{omfigure}

\begin{remark}[गुरुत्वाकर्षण का हमशक्ल जुड़वाँ]\label{rem:g11:fundamental-interactions:twin}
\omterm{def:g11:fundamental-interactions:charge}{कूलॉम} का नियम, चिह्न-दर-चिह्न, सार्वत्रिक \omterm{def:g11:fundamental-interactions:four}{गुरुत्वाकर्षण} का नियम
(\cref{ch:g10:universal-gravitation}) ही है, बस द्रव्यमानों की जगह आवेश रख दिए गए हैं:
\begin{center}
\begin{tabular}{lll}
स्रोत & द्रव्यमान $m$ & आवेश $q$ \\
नियम & $F = G\,\dfrac{m_1 m_2}{d^2}$ & $F = k\,\dfrac{\abs{q_1 q_2}}{d^2}$ \\[1.5ex]
अचर & $G = \qty{6.67e-11}{N.m^2/kg^2}$ & $k = \qty{8.99e9}{N.m^2/C^2}$ \\
चिह्न & सदा आकर्षी & आकर्षी या प्रतिकर्षी \\
\end{tabular}
\end{center}
नीचे की हर बात इन्हीं दो अंतरों से निकलती है: बिजली कहीं अधिक प्रबल है, और
केवल वही ख़ुद को काट भी सकती है।
\end{remark}

\begin{example}[दो प्रोटॉन, दोनों बल एक साथ]\label{ex:g11:fundamental-interactions:protons}
दो प्रोटॉन ($m_p = \qty{1.67e-27}{kg}$) $d = \qty{1.0e-15}{m}$ की दूरी पर बैठे हैं --- नाभिक के पड़ोसी। विद्युत
प्रतिकर्षण:
\[
F_E = \num{8.99e9} \times
\frac{(\num{1.60e-19})^{2}}{(\num{1.0e-15})^{2}}
\approx \qty{2.3e2}{N}
\]
--- यानी $23$ \omterm{def:g10:orders-of-magnitude:unit}{किलोग्राम} के सूटकेस का भार, और वह भी $10^{-27}$
\omterm{def:g10:orders-of-magnitude:unit}{किलोग्राम} के कण पर। गुरुत्वीय आकर्षण: $F_G = \num{6.67e-11} \times (\num{1.67e-27})^{2} /
(\num{1.0e-15})^{2} \approx \qty{1.9e-34}{N}$। अनुपात $F_E / F_G = k e^{2} / (G m_p^{2}) \approx \num{1.2e36}$ दूरी से
स्वतंत्र है ($d^{2}$ कट जाता है): मूल कणों के बीच गुरुत्व $10^{36}$ गुना कमज़ोर
है, इसलिए \emph{किसी भी} दूरी पर वह गिनती में नहीं आता।
\end{example}

\section{प्रबल और दुर्बल अंतःक्रिया}

उदाहरण एक कांड छोड़ जाता है: हीलियम के नाभिक में दो प्रोटॉन हैं जो
दसियों \omterm{def:g10:inertia:force}{न्यूटन} से प्रतिकर्षित करते हैं --- फिर भी हीलियम मौजूद है।

\begin{definition}[प्रबल अंतःक्रिया]\label{def:g11:fundamental-interactions:strong}
\emph{प्रबल अंतःक्रिया}\index{प्रबल अंतःक्रिया} क्वार्कों को बाँधकर
\omterm{def:g11:fundamental-interactions:ladder}{न्यूक्लियॉन} बनाती है और न्यूक्लियॉनों को बाँधकर नाभिक।
\qty{e-15}{m} पर यह आकर्षी है और विद्युत प्रतिकर्षण से कहीं अधिक प्रबल
(न्यूक्लियॉनों के बीच हज़ारों \omterm{def:g10:inertia:force}{न्यूटन}); पर चंद $10^{-15}$
\omterm{def:g10:orders-of-magnitude:unit}{मीटर} से आगे वह लुप्त हो जाती है: उसका
\emph{परास}\index{परास (अंतःक्रिया का)} लगभग \qty{e-15}{m} है। वह प्रोटॉन और
न्यूट्रॉन दोनों को एक-सा जकड़ती है और आवेश की परवाह नहीं करती।
\end{definition}

\begin{omfigure}
\begin{tikzpicture}[font=\small]
  \draw[omDef, thick, fill=omDef!20] (0,0) circle (0.34) node[black] {p};
  \draw[omDef, thick, fill=omDef!20] (3.2,0) circle (0.34) node[black] {p};
  \draw[-Stealth, omThm, thick] (-0.34,0) -- (-1.7,0)
    node[midway, above] {$F_E \approx \qty{230}{N}$};
  \draw[-Stealth, omThm, thick] (3.54,0) -- (4.9,0) node[midway, above] {$F_E$};
  \draw[-Stealth, omMeth, very thick] (0.2,0.55) -- (1.45,0.55)
    node[midway, above] {प्रबल};
  \draw[-Stealth, omMeth, very thick] (3.0,0.55) -- (1.75,0.55)
    node[midway, above] {प्रबल};
  \draw[Stealth-Stealth] (0,-0.65) -- (3.2,-0.65)
    node[midway, below] {$d = \qty{1.0e-15}{m}$};
\end{tikzpicture}

{\small नाभिक के भीतर की रस्साकशी: $10^{-15}$~m पर प्रबल आकर्षण (हज़ारों
\omterm{def:g10:inertia:force}{न्यूटन}) विद्युत प्रतिकर्षण को हरा देता है।}
\end{omfigure}

\begin{remark}[नाभिक क्यों टिकते हैं --- और भारी नाभिक क्यों हार जाते हैं]\label{rem:g11:fundamental-interactions:nuclei}
नाभिक के भीतर दोनों बल अलग-अलग ढंग से बढ़ते हैं। प्रबल गोंद
कम \omterm{def:g11:fundamental-interactions:strong}{परास} वाली है: एक \omterm{def:g11:fundamental-interactions:ladder}{न्यूक्लियॉन} \qty{e-15}{m} के भीतर बैठे अपने चंद
पड़ोसियों से ही बँधता है। विद्युत प्रतिकर्षण लंबी \omterm{def:g11:fundamental-interactions:strong}{परास} वाला है: हर प्रोटॉन
\emph{हर दूसरे} प्रोटॉन को धकेलता है। नाभिक जैसे-जैसे बढ़ते हैं, प्रतिकर्षण
गोंद से तेज़ जमा होता जाता है: भारी नाभिकों को अतिरिक्त न्यूट्रॉन चाहिए (आवेश
के बिना गोंद), और लगभग $80$ प्रोटॉन से आगे वह भी काम नहीं आता --- सबसे
भारी नाभिक उधार के समय पर जीते हैं (\cref{ch:g11:nucleus-radioactivity})।
\end{remark}

\begin{definition}[दुर्बल अंतःक्रिया]\label{def:g11:fundamental-interactions:weak}
\emph{दुर्बल अंतःक्रिया}\index{दुर्बल अंतःक्रिया}, जिसका \omterm{def:g11:fundamental-interactions:strong}{परास} और भी
छोटा है, न्यूट्रॉन को प्रोटॉन में बदल देने देती है (और उलटा भी): वही
\cref{ch:g11:nucleus-radioactivity} की $\beta$ रेडियोसक्रियता चलाती है।
\end{definition}

\section{कौन किस मंज़िल का राजा है}

\begin{method}[किसी पैमाने का लेखा-जोखा]\label{met:g11:fundamental-interactions:audit}
यह जानने के लिए कि $d$ आकार की किसी संरचना पर कौन-सी अंतःक्रिया राज करती
है, \omterm{def:g11:fundamental-interactions:strong}{परास} और कटाव, दोनों देखिए:
\begin{enumerate}
  \item $d \leq \qty{e-14}{m}$: \emph{\omterm{def:g11:fundamental-interactions:strong}{प्रबल अंतःक्रिया}} राज करती है --- वह विद्युत
        प्रतिकर्षण को हरा देती है, और गुरुत्व $10^{36}$ बाहर है।
  \item परमाणुओं से लेकर पहाड़ों तक ($\qty{e-10}{m}$ से $\qty{e4}{m}$ तक): प्रबल
        बल \omterm{def:g11:fundamental-interactions:strong}{परास} से बाहर है; \emph{\omterm{def:g11:fundamental-interactions:four}{विद्युत्-चुंबकीय
        अंतःक्रिया}} राज करती है --- बंध, संसंजन, \omterm{def:g10:inertia:inventory}{घर्षण}, संस्पर्श।
  \item ग्रह और उससे आगे ($d \geq \qty{e6}{m}$): पदार्थ अद्भुत परिशुद्धता से उदासीन है,
        इसलिए बिजली ख़ुद को काट देती है; द्रव्यमान का एक ही चिह्न होता है और
        वह सदा जुड़ता जाता है --- आसमान पर \emph{\omterm{def:g11:fundamental-interactions:four}{गुरुत्वाकर्षण}} का
        राज है।
\end{enumerate}
\omterm{def:g11:fundamental-interactions:weak}{दुर्बल अंतःक्रिया} के पास कोई मंज़िल नहीं है: वह कुछ नहीं बाँधती, पर
रूपांतरणों का फ़ैसला करती है ($\beta$ क्षय)।
\end{method}

\begin{omfigure}
\begin{tikzpicture}[font=\small, x=0.26cm]
  \draw[-Stealth] (-1,0) -- (41.5,0);
  \foreach \x/\l in {3/{10^{-15}}, 8/{10^{-10}}, 18/{10^{0}},
                     24/{10^{6}}, 39/{10^{21}}}
    \draw (\x,0.1) -- (\x,-0.1) node[below] {$\l$};
  \draw[omProp, fill=omProp!35] (3,0.3) rectangle (4,0.75);
  \node[anchor=west, omProp] at (4.4,0.52) {प्रबल};
  \draw[omDef, fill=omDef!25] (8,0.95) rectangle (22,1.4);
  \node[omDef] at (15,1.17) {विद्युत्-चुंबकीय};
  \draw[omThm, fill=omThm!25] (24,1.6) rectangle (39,2.05);
  \node[omThm] at (31.5,1.82) {गुरुत्वाकर्षण};
\end{tikzpicture}

{\small कहाँ किसका राज: नाभिकों पर प्रबल, परमाणुओं से पहाड़ों तक
विद्युत्-चुंबकीय, ग्रहों और उससे आगे \omterm{def:g11:fundamental-interactions:four}{गुरुत्वाकर्षण}। आकार
\omterm{def:g10:orders-of-magnitude:unit}{मीटर} में।}
\end{omfigure}

\section{अभ्यास}

\begin{exercise}[$\star$]\label{exo:g11:fundamental-interactions:1}
हर आकार की \omterm{def:g10:orders-of-magnitude:oom}{परिमाण की कोटि} और सीढ़ी पर उसकी मंज़िल बताइए: प्रोटॉन
\qty{1.7e-15}{m}; पानी का \omterm{def:g11:fundamental-interactions:ladder}{अणु} \qty{2.8e-10}{m}; रेत का कण \qty{2e-4}{m}; पृथ्वी \qty{1.3e7}{m};
आकाशगंगा \qty{9.5e20}{m}।
\end{exercise}

\begin{exercise}[$\star$]\label{exo:g11:fundamental-interactions:2}
बालों पर रगड़ा गुब्बारा $q = \qty{-4.8e-9}{C}$ रखता है। उसने इलेक्ट्रॉन पाए हैं या खोए, और
कितने? इसके बाद बालों का आवेश क्या है, और कौन-सा नियम यह कहता है?
\end{exercise}

\begin{exercise}[$\star$]\label{exo:g11:fundamental-interactions:3}
आवेश $q_1 = \qty{+2.0e-6}{C}$ और $q_2 = \qty{-3.0e-6}{C}$ \qty{30}{cm} की दूरी पर हैं। बल निकालिए; आकर्षी
या प्रतिकर्षी? $q_1$ पर और $q_2$ पर लगते बलों की तुलना कीजिए।
\end{exercise}

\begin{exercise}[$\star$]\label{exo:g11:fundamental-interactions:4}
दो आवेश $F_0$ बल से आकर्षित करते हैं। यह क्या हो जाएगा यदि: दूरी
दुगुनी कर दी जाए; दूरी को $3$ से भाग दिया जाए; एक आवेश दुगुना कर दिया
जाए; दोनों आवेश \emph{और} दूरी, तीनों दुगुने कर दिए जाएँ?
\end{exercise}

\begin{exercise}[$\star$]\label{exo:g11:fundamental-interactions:5}
इनके लिए ज़िम्मेदार अंतःक्रिया का नाम लिखिए: (क) चंद्रमा का पृथ्वी की
परिक्रमा; (ख) नमक के क्रिस्टल का संसंजन; (ग) हीलियम के नाभिक का
संसंजन; (घ) कार्बन-14 का $\beta$ क्षय; (ङ) \cref{exo:g11:fundamental-interactions:2} वाले गुब्बारे का दीवार से
चिपकना।
\end{exercise}

\begin{exercise}[$\star\star$]\label{exo:g11:fundamental-interactions:6}
हीलियम के नाभिक के दोनों प्रोटॉन $d = \qty{2.0e-15}{m}$ की दूरी पर हैं ($m_p = \qty{1.67e-27}{kg}$)।
उनका विद्युत प्रतिकर्षण, उनका गुरुत्वीय आकर्षण और दोनों का अनुपात निकालिए।
नाभिक को साबुत कौन रखता है?
\end{exercise}

\begin{exercise}[$\star\star$]\label{exo:g11:fundamental-interactions:7}
हाइड्रोजन के \omterm{def:g11:fundamental-interactions:ladder}{परमाणु} में इलेक्ट्रॉन ($m_e = \qty{9.11e-31}{kg}$) प्रोटॉन ($m_p = \qty{1.67e-27}{kg}$) से
$d = \qty{5.3e-11}{m}$ की दूरी पर है। उनके बीच का विद्युत बल और \omterm{def:g10:universal-gravitation:force}{गुरुत्वाकर्षण बल}
तथा उनका अनुपात निकालिए। परमाणुओं को कौन जोड़े रखता है?
\end{exercise}

\begin{exercise}[$\star\star$]\label{exo:g11:fundamental-interactions:8}
गोला A पर $q_A = \qty{+8.0e-9}{C}$ है; उसके जैसा ही गोला B उदासीन है। दोनों को छुआकर अलग करके
\qty{10}{cm} की दूरी पर रख दिया जाता है। हर एक पर कितना आवेश है (आवेश का संरक्षण
जाँचिए)? उनके बीच का बल निकालिए --- आकर्षी या प्रतिकर्षी?
\end{exercise}

\begin{exercise}[$\star\star$]\label{exo:g11:fundamental-interactions:9}
दो एक जैसे आवेश $q = \qty{1.0e-6}{C}$ \qty{0.90}{N} से प्रतिकर्षित करते हैं: वे कितनी दूर हैं?
बल गिरकर \qty{0.10}{N} कहाँ हो जाता है?
\end{exercise}

\begin{exercise}[$\star\star$]\label{exo:g11:fundamental-interactions:10}
पड़ोसी अणुओं के दो प्रोटॉन $d = \qty{1.0e-10}{m}$ की दूरी पर हैं। उनका विद्युत
प्रतिकर्षण निकालिए। इस दूरी पर \omterm{def:g11:fundamental-interactions:strong}{प्रबल अंतःक्रिया} का क्या योगदान है?
अणुओं को कौन-सी अंतःक्रिया बाँधती है?
\end{exercise}

\begin{exercise}[$\star\star$]\label{exo:g11:fundamental-interactions:11}
दो सिक्कों में से हर एक में लगभग $\num{e23}$ इलेक्ट्रॉन हैं। उनके बीच
\emph{दस लाख} में से एक इलेक्ट्रॉन सरका दीजिए और सिक्कों को \qty{1.0}{m} दूर रख
दीजिए। हर सिक्के का आवेश निकालिए, फिर बल। यह कितने द्रव्यमान के
भार के बराबर है? रोज़मर्रा के पदार्थ की उदासीनता पर निष्कर्ष
लिखिए।
\end{exercise}

\begin{exercise}[$\star\star\star$]\label{exo:g11:fundamental-interactions:12}
यूरेनियम के नाभिक का व्यास \qty{1.4e-14}{m} है और उसमें $92$ प्रोटॉन हैं।
आमने-सामने के दो सिरों पर बैठे दो प्रोटॉनों के बीच प्रतिकर्षण निकालिए। क्या
\omterm{def:g11:fundamental-interactions:strong}{प्रबल अंतःक्रिया} इस युग्म को बाँध सकती है --- क्यों नहीं? कुल कितने
प्रतिकर्षी प्रोटॉन-युग्म हैं? भारी नाभिकों को अतिरिक्त न्यूट्रॉन क्यों चाहिए,
और किसी आकार से आगे न्यूट्रॉन भी उन्हें क्यों नहीं बचा पाते?
\end{exercise}

\begin{exercise}[$\star\star\star$]\label{exo:g11:fundamental-interactions:13}
\qty{1.0}{g} द्रव्यमान की दो गेंदें एक ही बिंदु से बँधे \qty{50}{cm} के धागों पर लटकी
हैं। दोनों को एक ही आवेश $q$ देने पर वे \qty{6.0}{cm} की दूरी पर ठहर जाती हैं।
हर गेंद अपने भार, धागे के \omterm{def:g10:inertia:inventory}{तनाव} और विद्युत बल के
नीचे साम्यावस्था में है: दिखाइए कि $F = mg\tan\theta$ ($\theta$: धागे का \omterm{def:g10:universal-gravitation:weight}{ऊर्ध्वाधर} से
कोण) और यहाँ $\tan\theta \approx 0.060$। $F$ निकालिए, फिर $q$, फिर सरकाए गए
मूल आवेशों की संख्या।
\end{exercise}

\begin{exercise}[$\star\star\star$]\label{exo:g11:fundamental-interactions:14}
आँकड़े: $M_E = \qty{5.97e24}{kg}$, $M_M = \qty{7.35e22}{kg}$, पृथ्वी--चंद्रमा दूरी $d = \qty{3.84e8}{m}$। पृथ्वी--चंद्रमा
\omterm{def:g10:universal-gravitation:force}{गुरुत्वाकर्षण बल} निकालिए। मान लीजिए गुरुत्व बंद कर दिया गया: पृथ्वी
पर कितना आवेश $+Q$ और चंद्रमा पर कितना $-Q$ रखने से वही आकर्षण बना
रहेगा? $Q$ कितने इलेक्ट्रॉन हैं, और उनका भार कितना है ($m_e = \qty{9.11e-31}{kg}$)? टिप्पणी
कीजिए।
\end{exercise}

\begin{exercise}[$\star\star\star$]\label{exo:g11:fundamental-interactions:15}
नमक के क्रिस्टल में Na$^{+}$ और Cl$^{-}$ आयन (आवेश $\pm e$) $d = \qty{2.8e-10}{m}$ की
दूरी पर हैं; Cl$^{-}$ आयन का द्रव्यमान \qty{5.9e-26}{kg} है। पड़ोसी आयनों के बीच का
बल और Cl$^{-}$ आयन का भार निकालिए; अनुपात दीजिए। फिर
समझाइए कि हर बंध के गुरुत्व से पंद्रह कोटि आगे होने के बावजूद बड़े पैमानों पर
गुरुत्व कैसे जीत जाता है --- कोई पहाड़ \qty{e4}{m} से ऊँचा नहीं होता।
\end{exercise}

\section{समस्या: चार बलों का लेखा-जोखा}

\begin{problem}[{सप्ताहांत समस्या --- दुनिया न भरभराकर गिरती है न बिखर जाती
है, ऐसा क्यों: चारों अंतःक्रियाओं का मंज़िल-दर-मंज़िल लेखा, और पदार्थ की
$10^{18}$ में दो भाग तक की उदासीनता}]
\label{pb:g11:fundamental-interactions:1}
\omterm{def:g11:fundamental-interactions:ladder}{परमाणु} भीतर की ओर नहीं ढहते, नाभिक अधिकतर टिके रहते हैं, चंद्रमा न
गिरता है न भाग निकलता है; यह समस्या हिसाब लगाती है कि श्रेय किसे मिलना चाहिए।
आँकड़े: $e = \qty{1.60e-19}{C}$, $k = \qty{8.99e9}{N.m^2/C^2}$, $G = \qty{6.67e-11}{N.m^2/kg^2}$, $m_p = \qty{1.67e-27}{kg}$, $m_e = \qty{9.11e-31}{kg}$, $M_E = \qty{5.97e24}{kg}$, $M_M = \qty{7.35e22}{kg}$,
पृथ्वी--चंद्रमा दूरी \qty{3.84e8}{m}।

\medskip
\noindent\textbf{भाग I --- सीढ़ी।}
\begin{enumerate}
  \item \omterm{def:g11:fundamental-interactions:ladder}{क्वार्क} से आकाशगंगा तक सीढ़ी की मंज़िलें गिनाइए, हर एक के
        आकार की एक \omterm{def:g10:orders-of-magnitude:oom}{परिमाण की कोटि} के साथ।
  \item पैमाना-प्रतिरूप: नाभिक को \qty{1}{cm} का कंचा बना दीजिए; तब
        \omterm{def:g11:fundamental-interactions:ladder}{परमाणु} कितना चौड़ा होगा?
  \item \omterm{def:g11:fundamental-interactions:ladder}{परमाणु} के आयतन का कितना अंश उसका नाभिक घेरता है?
  \item एक बाल \qty{70}{\micro m} मोटा है। उसकी चौड़ाई में कितने \omterm{def:g11:fundamental-interactions:ladder}{परमाणु} आएँगे?
  \item \omterm{def:g11:fundamental-interactions:ladder}{क्वार्क} की छत (\qty{e-18}{m}) से आकाशगंगा (\qty{e21}{m}) तक दस की कितनी
        घातें हैं?
\end{enumerate}

\medskip
\noindent\textbf{भाग II --- विद्युत मंज़िल।}
\begin{enumerate}[resume]
  \item हाइड्रोजन के \omterm{def:g11:fundamental-interactions:ladder}{परमाणु} में इलेक्ट्रॉन प्रोटॉन से $d = \qty{5.3e-11}{m}$ की
        दूरी पर है: उनके बीच का विद्युत बल निकालिए।
  \item उनके बीच का \omterm{def:g10:universal-gravitation:force}{गुरुत्वाकर्षण बल} और दोनों का अनुपात निकालिए।
        \omterm{def:g11:fundamental-interactions:ladder}{परमाणु} को इनमें से कौन जोड़े रखता है?
  \item यदि परमाणुओं के भीतर गुरुत्व बंद कर दिया जाए तो क्या बदलेगा?
        और यदि बिजली बंद कर दी जाए?
  \item रोज़मर्रा के ``संस्पर्श'' बल --- फ़र्श का आपके पैरों को
        धकेलना, \omterm{def:g10:inertia:inventory}{घर्षण} --- असल में भेस बदले विद्युत्-चुंबकीय बल ही
        क्यों हैं?
  \item दो सिक्कों में से हर एक में लगभग $\num{e23}$ इलेक्ट्रॉन हैं। उनके बीच
        \emph{एक अरब} में से एक इलेक्ट्रॉन सरकाकर उन्हें \qty{1.0}{m} दूर रखिए:
        आवेश और बल निकालिए। निष्कर्ष?
  \item \qty{e-10}{m} से \qty{e4}{m} तक की मंज़िलों पर किसका राज है, और
        \omterm{def:g11:fundamental-interactions:strong}{प्रबल अंतःक्रिया} मुक़ाबले में क्यों नहीं है?
\end{enumerate}

\medskip
\noindent\textbf{भाग III --- नाभिकीय मंज़िल।}
\begin{enumerate}[resume]
  \item किसी नाभिक में दो प्रोटॉन \qty{1.0e-15}{m} की दूरी पर हैं: उनका
        विद्युत प्रतिकर्षण निकालिए।
  \item बिना किसी रोक के यह बल एक प्रोटॉन को कितना त्वरण देगा?
        $g$ से तुलना कीजिए।
  \item \omterm{def:g11:fundamental-interactions:strong}{प्रबल अंतःक्रिया} में वे तीन गुण गिनाइए जो यह समझाने के लिए
        चाहिए कि नाभिक क्यों टिकते हैं, न्यूट्रॉन भी क्यों जकड़े जाते हैं और
        नाभिक इतने छोटे क्यों हैं।
  \item यूरेनियम के नाभिक (व्यास \qty{1.4e-14}{m}, $92$ प्रोटॉन) में
        आमने-सामने के दो प्रोटॉनों के बीच प्रतिकर्षण निकालिए। क्या
        \omterm{def:g11:fundamental-interactions:strong}{प्रबल अंतःक्रिया} उस युग्म को सीधे बाँध सकती है? समझाइए कि
        नाभिक बढ़ने के साथ प्रतिकर्षण क्यों जीतने लगता है, और न्यूट्रॉन किस
        काम आते हैं।
  \item जैसा वादा था, एक वाक्य में: \omterm{def:g11:fundamental-interactions:weak}{दुर्बल अंतःक्रिया} करती क्या है?
\end{enumerate}

\medskip
\noindent\textbf{भाग IV --- खगोलीय मंज़िल, और फ़ैसला।}
\begin{enumerate}[resume]
  \item पृथ्वी और चंद्रमा के बीच का \omterm{def:g10:universal-gravitation:force}{गुरुत्वाकर्षण बल} निकालिए।
  \item पृथ्वी में लगभग $\num{1.8e51}$ प्रोटॉन हैं (और उतने ही इलेक्ट्रॉन), चंद्रमा
        में लगभग $\num{2.2e49}$। हर पिंड के प्रोटॉन-आवेश का कितना अंश $f$ बिना
        कटे छूट जाए तो विद्युत बल गुरुत्वीय बल के बराबर हो जाएगा?
  \item प्रश्न 7 का कमज़ोर \omterm{def:g11:fundamental-interactions:four}{गुरुत्वाकर्षण} खगोल पर राज क्यों करता है?
        दो कारण।
  \item फ़ैसला --- हर मंज़िल पर एक वाक्य (नाभिक, \omterm{def:g11:fundamental-interactions:ladder}{परमाणु} से
        पहाड़ तक, ग्रह और उससे ऊपर): उसे कौन जोड़े रखता है? फिर शीर्षक का
        उत्तर दीजिए: दुनिया न भरभराकर गिरती है न बिखर जाती है, ऐसा क्यों?
\end{enumerate}
\end{problem}
